\begin{abstract}
Inspired by physical balance properties of weighted grids, this paper presents a rigorous geometric and topological analysis of magic squares through manifold learning techniques. We develop a novel embedding framework that maps magic squares to point clouds in $\mathbb{R}^3$, revealing their intrinsic geometric structure and connecting their mathematical properties to physical interpretations. Our analysis demonstrates that the space of magic squares exhibits a stratified manifold structure with two distinct geometric regimes, characterized by different local dimensions and topological features. Using tools from computational topology and stratification learning theory, we identify persistent homological features that distinguish magic squares from both random and structured null models, while providing rigorous algorithmic complexity bounds. Through our dual null model comparison framework, we establish that magic squares possess significant higher-dimensional structure beyond what would be expected from their basic constraints. Our methodology combines Earth Mover's Distance metrics, dimension estimation algorithms, and persistent homology to create a comprehensive geometric theory of these classical mathematical objects. The analysis reveals a rich multi-scale organization with hierarchical component structure, significant cyclic features, and higher-dimensional void structures that provide new insights into the geometric constraints governing magic squares. This approach opens new avenues for analyzing other combinatorial structures through the lens of geometry and topology, while suggesting potential applications in physical systems, algorithm design, and discrete optimization.
\end{abstract} 